% !TEX program = pdflatex
% Hugging Face Deep RL Course -- Bonus Unit 5. Полный исходный материал MDX.
% Компиляция из notes/ru/: latexmk -pdf unitbonus5ru.tex
\PassOptionsToPackage{T2A}{fontenc}
\PassOptionsToPackage{russian}{babel}
\documentclass[11pt,a4paper]{article}

% Shared typography/layout for Hugging Face Deep RL Course notes.
% Compile with pdfLaTeX / latexmk -pdf with notes/ as the working directory.
% e.g. cd notes && latexmk -pdf unit1.tex

\usepackage[T1]{fontenc}
\usepackage[utf8]{inputenc}
\usepackage[english]{babel}
\usepackage{lmodern}
\usepackage{textcomp}
\usepackage{geometry}
\geometry{a4paper,margin=22mm,headheight=14pt,headsep=8mm,footskip=12mm}

\usepackage{amsmath,amssymb,mathtools}
\usepackage{graphicx,caption,needspace,placeins}
\usepackage{array,booktabs,tabularx}
\usepackage{xcolor}
\usepackage{microtype}
\usepackage{enumitem}
\usepackage{listings}
\usepackage[most]{tcolorbox}
\usepackage{titlesec}
\usepackage{fancyhdr}
\usepackage{hyperref}
\usepackage{bookmark}

\definecolor{ink}{HTML}{243041}
\definecolor{accent}{HTML}{38598C}
\definecolor{secondary}{HTML}{60768A}
\definecolor{codebg}{HTML}{F4F6F8}
\definecolor{codeborder}{HTML}{E3E8EE}
\definecolor{notebg}{HTML}{F3F7FC}

\hypersetup{
  colorlinks=true,
  linkcolor=accent,
  urlcolor=accent,
  citecolor=accent,
  pdfborder={0 0 0}
}
\urlstyle{same}

\setlength{\parindent}{0pt}
\setlength{\parskip}{0.38em plus 0.12em minus 0.05em}

\setlist[itemize]{
  topsep=0.35em,
  itemsep=0.13em,
  leftmargin=1.7em
}
\setlist[enumerate]{
  topsep=0.35em,
  itemsep=0.13em,
  leftmargin=1.8em
}

\setcounter{tocdepth}{2}
\setcounter{secnumdepth}{2}

\titleformat{\section}
  {\Large\bfseries\color{ink}}
  {\thesection}{0.7em}{}

\titleformat{\subsection}
  {\large\bfseries\color{accent}}
  {\thesubsection}{0.7em}{}

\titleformat{\subsubsection}
  {\normalsize\bfseries\color{ink}}
  {\thesubsubsection}{0.6em}{}

\titlespacing*{\section}
  {0pt}{2.2ex plus 0.7ex}{0.9ex}
\titlespacing*{\subsection}
  {0pt}{1.7ex plus 0.4ex}{0.6ex}
\titlespacing*{\subsubsection}
  {0pt}{1.3ex plus 0.3ex}{0.45ex}

\pagestyle{fancy}
\fancyhf{}

\fancyhead[L]{
  \footnotesize\sffamily\color{secondary}
  HF Deep RL Course \enspace/\enspace Unit 1
}
\fancyhead[R]{
  \footnotesize\sffamily\color{secondary}
  \nouppercase{\rightmark}
}
\fancyfoot[C]{
  \footnotesize\color{secondary}\thepage
}
\renewcommand{\headrulewidth}{0.25pt}

\captionsetup{
  font=small,
  labelfont={bf,color=secondary},
  textfont={color=ink},
  justification=centering,
  skip=4pt,
  hypcap=false
}

\lstdefinestyle{hfcode}{
  basicstyle=\ttfamily\footnotesize,
  backgroundcolor=\color{codebg},
  frame=single,
  rulecolor=\color{codeborder},
  numbers=none,
  showstringspaces=false,
  breaklines=true,
  breakatwhitespace=false,
  columns=fullflexible,
  keepspaces=true,
  xleftmargin=0pt,
  xrightmargin=0pt,
  aboveskip=0.6em,
  belowskip=0.6em,
  framexleftmargin=7pt,
  framexrightmargin=7pt,
  framextopmargin=5pt,
  framexbottommargin=5pt,
  upquote=true,
  tabsize=4,
}
\lstset{style=hfcode}

\newcommand{\CourseFigure}[2]{%
  \par\addvspace{0.55em}\noindent%
  \begin{minipage}{\linewidth}%
    \centering
    \includegraphics[
      width=0.86\linewidth,
      height=0.295\textheight,
      keepaspectratio
    ]{#1}\par
    \captionof{figure}{#2}%
  \end{minipage}%
  \par\addvspace{0.55em}%
}

\newtcolorbox{coursenote}{
  enhanced,
  breakable,
  colback=notebg,
  colframe=accent!30,
  boxrule=0.35pt,
  arc=1.2mm,
  left=9pt,
  right=9pt,
  top=6pt,
  bottom=6pt
}

\renewenvironment{quote}
  {\begin{coursenote}}
  {\end{coursenote}}

\emergencystretch=2.5em

\providecommand{\tightlist}{
  \setlength{\itemsep}{0pt}
  \setlength{\parskip}{0pt}
}

\providecommand{\passthrough}[1]{#1}

\Urlmuskip=0mu plus 6mu\relax

\renewcommand{\rmdefault}{cmr}
\renewcommand{\sfdefault}{cmss}
\renewcommand{\ttdefault}{cmtt}
\graphicspath{{../assets/unitbonus5/}}
\fancyhead[L]{\footnotesize\sffamily\color{secondary}Курс Hugging Face Deep RL / Бонусный модуль 5}
\fancyhead[R]{\footnotesize\sffamily\color{secondary}Обучение по демонстрациям}
\title{\textbf{Глубокое обучение с подкреплением}\\[0.35em]\large Бонусный модуль 5: имитационное обучение в Godot}
\author{Курс Hugging Face Deep RL\\\small Полная русскоязычная учебная версия на основе MDX}
\date{}
\begin{document}
\selectlanguage{russian}
\fontencoding{T2A}\selectfont
\maketitle
\vspace{-0.65em}
\begin{center}\small\color{secondary}Источник: курс Hugging Face Deep RL (перевод полного английского текста).\end{center}
\begingroup\small\setlength{\parskip}{0pt}\tableofcontents\endgroup
\vspace{0.4em}

\section{Введение}\label{introduction}

\CourseFigure{thumbnail.png}{Обучение по демонстрациям в Godot RL Agents: обзор модуля \protect\href{https://huggingface.co/datasets/huggingface-deep-rl-course/course-images/resolve/main/en/unit13/thumbnail.png}{[оригинал изображения]}}

Добро пожаловать в бонусный модуль, в котором вы \textbf{обучите агента-робота проходить уровень мини-игры с помощью имитационного обучения.}

К концу модуля \textbf{у вас будет обученный агент, способный проходить уровень так, как показано на видео}:

\begin{coursenote}\textbf{Видео:} обученный робот проходит уровень с помощью инференса ONNX\par\small\href{https://huggingface.co/datasets/huggingface-deep-rl-course/course-images/resolve/main/en/unit13/onnx_inference_test.mp4}{Посмотреть исходное видео}\end{coursenote}

\subsection{Цели}\label{objectives}

\begin{itemize}
\tightlist
\item
  Научиться применять имитационное обучение с Godot RL Agents: обучить агента проходить мини-игру на основе записанных человеком экспертных демонстраций.
\end{itemize}

\subsection{Предварительные знания и требования}\label{prerequisites-and-requirements}

\begin{itemize}
\tightlist
\item
  Перед началом рекомендуется пройти предыдущую главу (\href{https://huggingface.co/learn/deep-rl-course/unitbonus3/godotrl}{Godot RL Agents}).
\item
  Желательно некоторое знакомство с Godot, хотя для прохождения руководства писать код на GDScript не потребуется.
\item
  Godot с поддержкой .NET (работоспособность проверена с \href{https://godotengine.org/article/dev-snapshot-godot-4-3-dev-5/}{4.3.dev5 .NET}; более новые версии тоже могут подойти).
\item
  Godot RL Agents (в виртуальном окружении venv/conda можно выполнить \passthrough{\lstinline!pip install godot-rl!}).
\item
  \href{https://huggingface.co/learn/deep-rl-course/unitbonus5/train-our-robot}{Библиотека imitation}.
\item
  Время: приблизительно 1--2 часа на выполнение проекта и обучение. В зависимости от оборудования это может занять больше или меньше времени.
\end{itemize}

\section{Среда}\label{the-environment}

\CourseFigure{il_tutorial_env.png}{Уровень мини-игры: рычаг, ключ и сундук с сокровищами \protect\href{https://huggingface.co/datasets/huggingface-deep-rl-course/course-images/resolve/main/en/unit13/il_tutorial_env.png}{[оригинал изображения]}}

В этой учебной среде робот должен выполнить следующие действия:

\begin{itemize}
\tightlist
\item
  Потянуть за рычаг, чтобы поднять лестницу, ведущую во вторую комнату.
\item
  Добраться до ключа \texttt{[key]} и подобрать его, не упав в ловушки, воду или за пределы карты.
\item
  Вернуться к сундуку с сокровищами в первой комнате и открыть его. Победа! \texttt{[trophy]}
\end{itemize}

\section{Начало работы}\label{getting-started}

Для начала скачайте проект \href{https://huggingface.co/ivan267/imitation-learning-tutorial-godot-project/tree/main}{по ссылке} (нажмите значок загрузки рядом с файлом \passthrough{\lstinline!GDRL-IL-Project.zip!}). ZIP-архив содержит две версии проекта: «Starter» (начальная) и «Complete» (завершённая).

Игровая логика в начальном проекте уже реализована, а узлы настроены. Мы сосредоточимся на следующих задачах:

\begin{itemize}
\tightlist
\item
  Реализация кода узла AIController.
\item
  Запись экспертных демонстраций.
\item
  Обучение агента и экспорт файла .onnx для последующего инференса в Godot.
\end{itemize}

\subsection{Откройте начальный проект в Godot}\label{open-the-starter-project-in-godot}

Распакуйте ZIP-архив, откройте Godot, нажмите «Import» и укажите папку \passthrough{\lstinline!Starter\\Godot!} из распакованного архива.

\subsection{Откройте сцену робота}\label{open-the-robot-scene}

\begin{quote}
\textbf{Совет.} Сцену можно найти по слову «robot» через поиск FileSystem.
\end{quote}

В этой сцене есть несколько разных узлов. Узел \passthrough{\lstinline!robot!} содержит визуальную модель робота. Узел \passthrough{\lstinline!CameraXRotation!} отвечает за поворот камеры мышью вверх и вниз при ручном управлении. ИИ-агент этим узлом не управляет: для обучения решению задачи он не нужен. Узел \passthrough{\lstinline!RaycastSensors!} содержит два датчика Raycast, которые позволяют агенту «ощущать» части игрового мира: стены, пол и т. п.

\CourseFigure{open-robot-scene.jpg}{Редактор Godot: открытие сцены робота \protect\href{https://huggingface.co/datasets/huggingface-deep-rl-course/course-images/resolve/main/en/unit13/open-robot-scene.jpg}{[оригинал изображения]}}

\subsection{Нажмите значок скрипта рядом с AIController3D, чтобы открыть его для редактирования}\label{click-on-the-scroll-next-to-aicontroller3d-to-open-the-script-for-editing}

\begin{quote}
\textbf{Совет.} Чтобы проще найти узел, может потребоваться свернуть ветку «robot». Можно также ввести \passthrough{\lstinline!aicontroller!} в поле Filter над узлом \passthrough{\lstinline!Robot!}.
\end{quote}

\subsection{Замените методы \texttt{get\_obs()} и \texttt{get\_reward()} приведённой ниже реализацией}\label{replace-the-get_obs-and-get_reward-methods-with-the-implementation-below}

\begin{lstlisting}[language=Python]
func get_obs() -> Dictionary:
	var observations: Array[float] = []
	for raycast_sensor in raycast_sensors:
		observations.append_array(raycast_sensor.get_observation())

	var level_size = 16.0

	var chest_local = to_local(chest.global_position)
	var chest_direction = chest_local.normalized()
	var chest_distance = clampf(chest_local.length(), 0.0, level_size)
	
	var lever_local = to_local(lever.global_position)
	var lever_direction = lever_local.normalized()
	var lever_distance = clampf(lever_local.length(), 0.0, level_size)
		
	var key_local = to_local(key.global_position)
	var key_direction = key_local.normalized()
	var key_distance = clampf(key_local.length(), 0.0, level_size)
	
	var raft_local = to_local(raft.global_position)
	var raft_direction = raft_local.normalized()
	var raft_distance = clampf(raft_local.length(), 0.0, level_size)
	
	var player_speed = player.global_basis.inverse() * player.velocity.limit_length(5.0) / 5.0

	(
		observations
		.append_array(
			[
				chest_direction.x,
				chest_direction.y,
				chest_direction.z,
				chest_distance,
				lever_direction.x,
				lever_direction.y,
				lever_direction.z,
				lever_distance,
				key_direction.x,
				key_direction.y,
				key_direction.z,
				key_distance,
				raft_direction.x,
				raft_direction.y,
				raft_direction.z,
				raft_distance,
				raft.movement_direction_multiplier,
				float(player._is_lever_pulled),
				float(player._is_chest_opened),
				float(player._is_key_collected),
				float(player.is_on_floor()),
				player_speed.x,
				player_speed.y,
				player_speed.z,
			]
		)
	)
	return {"obs": observations}

func get_reward() -> float:
	return reward
\end{lstlisting}

В \passthrough{\lstinline!get_obs()!} мы сначала получаем наблюдения от двух датчиков Raycast, добавленных к узлу \passthrough{\lstinline!AIController3D!} в инспекторе, и включаем их в массив наблюдений. Затем вычисляем векторы относительного положения сундука, рычага, ключа и плота, разделяем их на направления и расстояния и также добавляем в наблюдения.

Кроме того, в наблюдения включается другая информация о состоянии игры:

\begin{itemize}
\tightlist
\item
  был ли потянут рычаг;
\item
  подобран ли ключ;
\item
  открыт ли сундук;
\item
  стоит ли персонаж на полу (это также определяет, может ли он прыгнуть);
\item
  нормализованная локальная скорость персонажа.
\end{itemize}

Булевы значения, например \passthrough{\lstinline!_is_lever_pulled!}, преобразуются в числа с плавающей точкой (0 или 1).

В методе \passthrough{\lstinline!get_reward()!} достаточно вернуть текущее вознаграждение.

\subsection{Замените методы \texttt{\_physics\_process()} и \texttt{reset()} приведённой ниже реализацией}\label{replace-the-_physics_process-and-reset-methods-with-the-implementation-below}

\begin{lstlisting}[language=Python]
func _physics_process(delta: float) -> void:
	# Reset on timeout, this is implemented in parent class to set needs_reset to true,
	# we are re-implementing here to call player.game_over() that handles the game reset.
	n_steps += 1
	if n_steps > reset_after:
		player.game_over()

	# In training or onnx inference modes, this method will be called by sync node with actions provided,
	# For expert demo recording mode, it will be called without any actions (as we set the actions based on human input),
	# For human control mode the method will not be called, so we call it here without any actions provided.
	if control_mode == ControlModes.HUMAN:
		set_action()

	# Reset the game faster if the lever is not pulled.
	steps_without_lever_pulled += 1
	if steps_without_lever_pulled > 200 and (not player._is_lever_pulled):
		player.game_over()

func reset():
	super.reset()
	steps_without_lever_pulled = 0
\end{lstlisting}

\subsection{Замените методы \texttt{get\_action\_space()}, \texttt{get\_action()} и \texttt{set\_action()} приведённой ниже реализацией}\label{replace-the-get_action_space-get_action-and-set_action-methods-with-the-implementation-below}

\begin{lstlisting}[language=Python]
# Defines the actions for the AI agent ("size": 2 means 2 floats for this action)
func get_action_space() -> Dictionary:
	return {
		"movement": {"size": 2, "action_type": "continuous"},
		"rotation": {"size": 1, "action_type": "continuous"},
		"jump": {"size": 1, "action_type": "continuous"},
		"use_action": {"size": 1, "action_type": "continuous"}
	}

# We return the action values in the same order as defined in get_action_space() (important), but all in one array
# For actions of size 1, we return 1 float in the array, for size 2, 2 floats in the array, etc.
# set_action is called just before get_action by the sync node, so we can read the newly set values
func get_action():
	return [
		# "movement" action values
		player.requested_movement.x,
		player.requested_movement.y,
		# "rotation" action value
		player.requested_rotation.x,
		# "jump" action value (-1 if not requested, 1 if requested)
		-1.0 + 2.0 * float(player.jump_requested),
		# "use_action" action value (-1 if not requested, 1 if requested)
		-1.0 + 2.0 * float(player.use_action_requested)
	]

# Here we set human control and AI control actions to the robot
func set_action(action = null) -> void:
	# If there's no action provided, it means that AI is not controlling the robot (human control),
	if not action:
		# Only rotate if the mouse has moved since the last set_action call
		if previous_mouse_movement == mouse_movement:
			mouse_movement = Vector2.ZERO

		player.requested_movement = Input.get_vector(
			"move_left", "move_right", "move_forward", "move_back"
		)
		player.requested_rotation = mouse_movement

		var use_action = Input.is_action_pressed("requested_action")
		var jump = Input.is_action_pressed("requested_jump")

		player.use_action_requested = use_action
		player.jump_requested = jump

		previous_mouse_movement = mouse_movement	
	else:
		# If there is action provided, we set the actions received from the AI agent 
		player.requested_movement = Vector2(action.movement[0], action.movement[1])
		# The agent only rotates the robot along the Y axis, no need to rotate the camera along X axis
		player.requested_rotation = Vector2(action.rotation[0], 0.0)
		player.jump_requested = bool(action.jump[0] > 0)
		player.use_action_requested = bool(action.use_action[0] > 0)
\end{lstlisting}

Метод \passthrough{\lstinline!get_action()!} (нужный только в режиме записи демонстраций) должен возвращать действия, которые агенту следует выполнять при встрече с тем же состоянием. Важно, чтобы значения лежали в допустимом диапазоне (\passthrough{\lstinline!-1.0 to 1.0!}), поэтому для булевых состояний используется выражение \passthrough{\lstinline!-1 + 2 * variable!}. Кроме того, действия должны идти в правильном порядке, заданном в \passthrough{\lstinline!get_action_space()!}.

В режиме записи демонстраций \passthrough{\lstinline!set_action()!} вызывается без аргумента действий: значения задаются по вводу человека. В режимах обучения и инференса этот метод получает аргумент \passthrough{\lstinline!action!}, содержащий значения всех действий от RL-модели. Поэтому оба случая обрабатываются конструкцией \passthrough{\lstinline!if/else!}.

Дополнительные пояснения приведены в комментариях к коду.

\subsection{Замените метод \texttt{\_input} приведённой ниже реализацией}\label{replace-the-_input-method-with-the-implementation-below}

\begin{lstlisting}[language=Python]
# Record mouse movement for human and demo_record modes
# We don't directly rotate in input to allow for frame skipping (action_repeat setting) which
# will also be applied to the AI agent in training/inference modes.
func _input(event):
	if not (heuristic == "human" or heuristic == "demo_record"):
		return

	if event is InputEventMouseMotion:
		var movement_scale: float = 0.005
		mouse_movement.y = clampf(event.relative.y * movement_scale, -1.0, 1.0)
		mouse_movement.x = clampf(event.relative.x * movement_scale, -1.0, 1.0)
\end{lstlisting}

Этот фрагмент кода записывает движение мыши в режимах ручного управления и записи демонстраций.

\textbf{Наконец, сохраните скрипт. Теперь можно переходить к следующему шагу.}

\subsection{Откройте сцену записи демонстраций и выберите узел AIController3D}\label{open-the-demo-record-scene-and-click-on-aicontroller3d-node}

\begin{quote}
\textbf{Совет.} В поиске FileSystem можно ввести «demo», а в поле фильтрации узлов сцены --- «aicontroller».
\end{quote}

\CourseFigure{demo_record_scene.jpg}{Редактор Godot: сцена записи экспертных демонстраций \protect\href{https://huggingface.co/datasets/huggingface-deep-rl-course/course-images/resolve/main/en/unit13/demo_record_scene.jpg}{[оригинал изображения]}}

Ничего менять не нужно: все параметры уже заданы. Но разберём настройки, которые потребовались бы для вашей собственной среды.

В сцене изменены настройки узла \passthrough{\lstinline!Level > Robot > AIController3D!}:

\begin{itemize}
\tightlist
\item
  Параметр \passthrough{\lstinline!Control Mode!} установлен в \passthrough{\lstinline!Record Expert Demos!}.
\item
  Заполнено поле \passthrough{\lstinline!Expert Demo Save Path!}.
\item
  Значение \passthrough{\lstinline!Action Repeat!} совпадает со значением узла \passthrough{\lstinline!Sync!} в сценах \passthrough{\lstinline!training_scene!} и \passthrough{\lstinline!onnx_inference_scene!}. Это значит, что каждое действие агента повторяется на протяжении трёх физических кадров. Настройка в \passthrough{\lstinline!AIController!} обеспечивает такое же повторение действий при управлении человеком (из-за чего возникает небольшая задержка), чтобы поведение соответствовало поведению агента. Значение довольно мало и не создаёт заметной задержки. Если решите его изменить, поменяйте его во всех трёх местах.
\item
  Настройка \passthrough{\lstinline!Remove Last Episode!} позволяет назначить клавишу для удаления неудачного эпизода прямо во время записи, без перезапуска сеанса. Например, если робот упал в воду и игра сбросилась, этой клавишей можно удалить предыдущий эпизод, уже записывая следующий. По умолчанию выбрана \passthrough{\lstinline!R!}, но назначение можно изменить: нажмите на него, затем на кнопку \passthrough{\lstinline!Configure!}.
\end{itemize}

Ещё один способ упростить запись эпизодов в сложной среде --- замедлить её во время записи. Для этого выберите узел \passthrough{\lstinline!Sync!} в сцене и измените свойство \passthrough{\lstinline!Speed Up!} (по умолчанию равно 1).

\subsection{Запишем несколько демонстраций}\label{lets-record-some-demos}

\begin{quote}
\textbf{Совет.} Демонстрации сохраняются, только если записан хотя бы один полный эпизод, а окно игры закрыто кнопкой «X» или сочетанием ALT+F4. Кнопка остановки в редакторе Godot не сохраняет демонстрации. Лучше сначала записать один эпизод и проверить, появился ли файл \passthrough{\lstinline!expert_demos.json!} в файловой системе или папке проекта Godot.
\end{quote}

Убедитесь, что открыта сцена \passthrough{\lstinline!demo_record_scene!}, и \passthrough{\lstinline!press F6!}, чтобы начать запись демонстраций.

Управление:

\begin{itemize}
\tightlist
\item
  мышь управляет камерой (если нужно изменить чувствительность, откройте сцену \passthrough{\lstinline!robot!}, выберите узел \passthrough{\lstinline!Robot!} и измените параметр \passthrough{\lstinline!Rotation Speed!}; сохраняйте одинаковое значение при записи, обучении и инференсе);
\item
  клавиши \passthrough{\lstinline!WASD!} отвечают за перемещение персонажа;
\item
  \passthrough{\lstinline!SPACE!} --- прыжок;
\item
  \passthrough{\lstinline!E!} --- активация рычага и открытие сундука.
\end{itemize}

Сначала можете немного потренироваться, чтобы привыкнуть к среде. Если хотите пропустить запись демонстраций, в завершённом проекте уже есть готовый файл \passthrough{\lstinline!expert_demos.json!}, которым можно воспользоваться.

Записанные демонстрации должны содержать не менее 22--24 полных успешных эпизодов. При обучении можно использовать несколько файлов демонстраций, поэтому необязательно записывать всё за один раз (имя файла задаётся упомянутым выше параметром \passthrough{\lstinline!Expert Demo Save Path!}).

Запись 23 эпизодов заняла у автора около \textasciitilde10 минут. Ключ появляется поочерёдно в двух местах, поэтому 22 или 24 эпизода обеспечили бы равномерное распределение его положений в демонстрациях, хотя и 23 достаточно близко. Подходя к рычагу или сундуку, автор удерживал клавишу \passthrough{\lstinline!E!} чуть дольше, чтобы действие попало в несколько шагов рядом с объектом. Несколько неудачных эпизодов он удалил клавишей \passthrough{\lstinline!R!} во время записи следующего эпизода.

Ниже --- ускоренное видео процесса записи демонстраций:

\begin{coursenote}\textbf{Видео:} запись экспертных демонстраций человеком (ускоренная съёмка)\par\small\href{https://huggingface.co/datasets/huggingface-deep-rl-course/course-images/resolve/main/en/unit13/demo_record.mp4}{Посмотреть исходное видео}\end{coursenote}

\subsection{Экспортируем игру для обучения}\label{export-the-game-for-training}

Экспортировать игру из Godot можно через меню \passthrough{\lstinline!Project > Export!}.

\section{Обучение нашего робота}\label{train-our-robot}

\begin{quote}
\textbf{Совет.} Прежде чем начать обучение, установите библиотеку \href{https://imitation.readthedocs.io/en/latest/getting-started/installation.html}{imitation} в то же виртуальное окружение venv или conda, в котором установлен Godot RL Agents. Для этого выполните \passthrough{\lstinline!pip install imitation!}.
\end{quote}

\subsection{Скачайте скрипт \href{https://github.com/edbeeching/godot_rl_agents/blob/main/examples/sb3_imitation.py}{имитационного обучения} из репозитория Godot RL}\label{download-a-copy-of-the-imitation-learning-script-from-the-godot-rl-repository.}

\subsection{Запустите обучение со следующими аргументами}\label{run-training-using-the-arguments-below}

\begin{lstlisting}[language=Python]
sb3_imitation.py --env_path="path_to_ILTutorial_executable" --bc_epochs=100 --gail_timesteps=1450000 --demo_files "path_to_expert_demos.json" --n_parallel=4 --speedup=20 --onnx_export_path=model.onnx --experiment_name=ILTutorial
\end{lstlisting}

\textbf{Укажите в параметре пути к среде путь до экспортированной игры, а в параметре файлов демонстраций --- путь к записанным данным. Если файлов несколько, перечислите их через пробел, например \passthrough{\lstinline!--demo_files demos.json demos2.json!}.}

Можно также задать большое число шагов в \passthrough{\lstinline!--gail_timesteps!}, а затем вручную остановить обучение сочетанием \passthrough{\lstinline!CTRL+C!}. Автор остановил обучение таким способом, когда вознаграждение стало приближаться к 3, на отметке \passthrough{\lstinline!total_timesteps | 1.38e+06!}. Это заняло примерно \textasciitilde41 минуту, а предварительное обучение BC --- около \textasciitilde5,5 минуты на CPU его компьютера.

Чтобы видеть среду во время обучения, добавьте аргумент \passthrough{\lstinline!--viz!}. На этапе BC среда будет неподвижна: для этого этапа она не используется, кроме получения информации о пространствах наблюдений и действий. На этапе GAIL визуализация среды начнёт обновляться.

Ниже приведены показатели \passthrough{\lstinline!ep_rew_mean!} и \passthrough{\lstinline!ep_rew_wrapped_mean!} из логов, отображённые в \href{https://github.com/edbeeching/godot_rl_agents/blob/main/docs/TRAINING_STATISTICS.md}{TensorBoard}. В данном случае их графики довольно близки:

\CourseFigure{training_results.png}{Кривые обучения: ep\_rew\_mean и ep\_rew\_wrapped\_mean в TensorBoard \protect\href{https://huggingface.co/datasets/huggingface-deep-rl-course/course-images/resolve/main/en/unit13/training_results.png}{[оригинал изображения]}}

\begin{quote}
\textbf{Совет.} Логи находятся в папке \passthrough{\lstinline!logs/ILTutorial!} относительно каталога, из которого запускалось обучение. Если выполняете несколько запусков, каждый раз меняйте аргумент \passthrough{\lstinline!--experiment_name!}.
\end{quote}

Хотя для самого обучения здесь не нужно задавать вознаграждения среды (и они не используются), для контроля успешности реализована простая разреженная функция вознаграждения. Падение за пределы карты, в воду или ловушки задаёт \passthrough{\lstinline!reward += -1!}, а активация рычага, подбор ключа и открытие сундука дают по \passthrough{\lstinline!reward += 1!}. Если \passthrough{\lstinline!ep_rew_mean!} приближается к 3, результат хороший. Показатель \passthrough{\lstinline!ep_rew_wrapped_mean!} --- вознаграждение от дискриминатора GAIL, которое напрямую не показывает, насколько успешно агент решает задачу.

\subsection{Проверим обученного агента}\label{lets-test-the-trained-agent}

После обучения в папке, из которой запускался скрипт, появится файл \passthrough{\lstinline!model.onnx!} (полный путь к файлу \passthrough{\lstinline!.onnx!} также можно найти ближе к концу лога обучения в консоли). \textbf{Скопируйте его в папку игрового проекта Godot.}

\subsection{Откройте сцену ONNX-инференса}\label{open-the-onnx-inference-scene}

Как и сцена записи демонстраций, эта сцена использует только один экземпляр уровня. Для её узла \passthrough{\lstinline!Sync!} задан режим \passthrough{\lstinline!Onnx Inference!}.

\textbf{Выберите узел \passthrough{\lstinline!Sync!} и задайте для свойства \passthrough{\lstinline!Onnx Model Path!} значение \passthrough{\lstinline!model.onnx!}.}

\CourseFigure{onnx_inference_scene.jpg}{Редактор Godot: настройка сцены ONNX-инференса \protect\href{https://huggingface.co/datasets/huggingface-deep-rl-course/course-images/resolve/main/en/unit13/onnx_inference_scene.jpg}{[оригинал изображения]}}

\textbf{Нажмите F6, чтобы запустить сцену, и посмотрим, чему научился агент!}

Видео работы обученного агента:

\begin{coursenote}\textbf{Видео:} обученный робот проходит уровень с помощью инференса ONNX\par\small\href{https://huggingface.co/datasets/huggingface-deep-rl-course/course-images/resolve/main/en/unit13/onnx_inference_test.mp4}{Посмотреть исходное видео}\end{coursenote}

Судя по результатам, агент умеет подбирать ключ в обоих возможных положениях (на левой или правой платформе) и хорошо воспроизводит записанное поведение. \textbf{Если получили похожий результат --- отлично, вы успешно прошли это руководство!} \texttt{[trophy]}\texttt{[applause]}

Если результаты существенно отличаются, учтите, что на них влияют количество и качество демонстраций. Может помочь изменение числа шагов на этапах BC/GAIL и гиперпараметров Python-скрипта. Есть и вариации от запуска к запуску, поэтому даже при одинаковых настройках результаты иногда немного различаются.

\section{(Необязательно) Как изменить среду}\label{optional-how-to-customize-the-environment}

Если хотите изменить игровой уровень, откройте сцену \passthrough{\lstinline!res://scenes/level.tscn!}, затем в FileSystem Godot откройте папку \passthrough{\lstinline!res://scenes/modules/!}:

\CourseFigure{level_scene.jpg}{Сцена уровня и переиспользуемые модули комнат в Godot \protect\href{https://huggingface.co/datasets/huggingface-deep-rl-course/course-images/resolve/main/en/unit13/level_scene.jpg}{[оригинал изображения]}}

Уровень содержит три комнаты, собранные из модулей, робота и дополнительные коллайдеры. Последние не позволяют пройти уровень обходным способом: забраться на стену первой комнаты и оттуда добраться до ключа. Добавляя модули в сцену, можно создавать новые комнаты и предметы.

Если выбрать узел Key (он находится в \passthrough{\lstinline!Room3!}, но его можно найти и поиском), а затем открыть \passthrough{\lstinline!Node > Signals!}, станет видно, что сигнал \passthrough{\lstinline!collected!} подключён и к роботу, и к сундуку. Это позволяет отслеживать подбор ключа и разблокировать сундук. Такой же механизм используется для активации лестницы рычагом. Если добавить новые рычаги, лестницы или ключи, их можно соединить сигналами.

\CourseFigure{level_signals.jpg}{Сигналы Godot: связывание собираемых предметов и событий уровня \protect\href{https://huggingface.co/datasets/huggingface-deep-rl-course/course-images/resolve/main/en/unit13/level_signals.jpg}{[оригинал изображения]}}

Если переключиться на \passthrough{\lstinline!Groups!}, видно, что ключ входит в группу \passthrough{\lstinline!resetable!}. В эту же группу входят плот, рычаг, сундук и персонаж. Сюда можно добавить любой узел, который требуется сбрасывать при перезапуске эпизода.

\CourseFigure{resetables_group.jpg}{Узлы группы resetable, сбрасываемые при перезапуске эпизода \protect\href{https://huggingface.co/datasets/huggingface-deep-rl-course/course-images/resolve/main/en/unit13/resetables_group.jpg}{[оригинал изображения]}}

Для корректной работы каждый объект группы \passthrough{\lstinline!resetable!} должен реализовать метод \passthrough{\lstinline!reset()!}, отвечающий за его сброс.

Поскольку во время обучения используется несколько экземпляров сцены уровня, мы сбрасываем не все объекты \passthrough{\lstinline!resetables!}, а только находящиеся внутри соответствующего экземпляра сцены. За это отвечает метод \passthrough{\lstinline!reset_all_resetables()!} в файле \passthrough{\lstinline!level_manager.gd!}. Скрипт робота вызывает его при необходимости сброса.

После изменения размера уровня необходимо обновить переменную \passthrough{\lstinline!level_size!} в файле \passthrough{\lstinline!robot_ai_controller.gd!}. Приблизительно измерьте наибольший размер уровня и запишите это значение в переменную.

Если изменяется количество объектов, которые должен отслеживать \passthrough{\lstinline!AIController!} (рычагов, плотов и т. п.), нужно соответствующим образом изменить код скрипта, добавить экспортируемые свойства для этих объектов, а затем привязать их через свойства инспектора \passthrough{\lstinline!AIController!} в сцене уровня:

\CourseFigure{ai_controller_inspector_properties.jpg}{Инспектор Godot: ссылки AIController на объекты уровня \protect\href{https://huggingface.co/datasets/huggingface-deep-rl-course/course-images/resolve/main/en/unit13/ai_controller_inspector_properties.jpg}{[оригинал изображения]}}

После этого, возможно, придётся обновить такие же свойства \passthrough{\lstinline!AIController!} в сцене записи демонстраций.

\section{Заключение}\label{conclusion}

\textbf{Поздравляем с завершением этого бонусного модуля!} Вы познакомились с процессом записи экспертных демонстраций и обучения агента методом имитационного обучения (IL). В некоторых задачах это может быть альтернативой обучению игровых агентов методами RL.

Автор руководства --- \href{https://github.com/Ivan-267}{Иван Додич (Ivan Dodic)}. Благодарим \href{https://twitter.com/edwardbeeching}{Эдварда Бичинга (Edward Beeching)} и \href{https://twitter.com/thomassimonini}{Томаса Симонини (Thomas Simonini)} за рецензирование и отзывы.

\end{document}
